% ================= SESSION 9.2 =================
\section{9.2 DQN의 안정화 장치와 실습}
\begin{frame}{두 불안정 원인: 이번 주차의 지도}
  \begin{columns}[T]
    \begin{column}{0.5\textwidth}
      \kb{문제 1: 목표가 움직인다}
      \begin{itemize}
        \item ``정답''인 목표값이 매 스텝 스스로 바뀜.
        \item 해결: \textbf{타겟 네트워크} (이 절).
      \end{itemize}
    \end{column}
    \begin{column}{0.5\textwidth}
      \kb{문제 2: 데이터가 상관된다}
      \begin{itemize}
        \item 연속 경험이 서로 비슷해서 미니배치 분산이 커짐.
        \item 해결: \textbf{경험 재현} (9.3절).
      \end{itemize}
    \end{column}
  \end{columns}
  \vskip0.6em
  \begin{block}{실험의 목표}
    ``움직이는 목표''가 학습을 \textbf{실제로 얼마나} 변질시키는지를
    숫자와 그래프로 확인하고, 그 실험이 건져 올린
    \textbf{세 번째 문제(과대추정)}의 실마리를 9.3절 Double DQN으로
    이어갈 고리를 남긴다.
  \end{block}
\end{frame}
